% !TEX program = xelatex
\documentclass[11pt,openany,fontset=fandol]{ctexbook}
\usepackage[paperwidth=176mm,paperheight=250mm,inner=21mm,outer=18mm,top=20mm,bottom=20mm,headheight=14pt,headsep=6mm,footskip=10mm]{geometry}
\usepackage{amsmath,amssymb,amsthm,bm,mathtools,booktabs,array,enumitem,float}
\usepackage{tikz,fancyhdr,hyperref}
\usetikzlibrary{arrows.meta,calc,patterns,angles,quotes,decorations.markings,decorations.pathmorphing}
\hypersetup{hidelinks}\setlength{\parindent}{2em}\setlength{\parskip}{3pt}\setlength{\emergencystretch}{2em}
\linespread{1.18}\raggedbottom\setcounter{tocdepth}{1}
\ctexset{chapter={name={第,讲},number=\chinese{chapter},format=\LARGE\bfseries\centering,beforeskip=16pt,afterskip=22pt},section={format=\Large\bfseries},subsection={format=\large\bfseries},subsubsection={format=\normalsize\bfseries,beforeskip=10pt,afterskip=4pt}}
\pagestyle{fancy}\fancyhf{}\fancyhead[LE,RO]{\small\thepage}\fancyhead[LO]{\small 物理竞赛一轮讲义}
\renewcommand{\headrulewidth}{0pt}\fancypagestyle{plain}{\fancyhf{}\fancyfoot[C]{\small\thepage}\renewcommand{\headrulewidth}{0pt}}
\theoremstyle{definition}\newtheorem{example}{例}[chapter]\newtheorem{drill}{练习}[chapter]
\widowpenalty=10000\clubpenalty=10000\displaywidowpenalty=10000
\newenvironment{solution}{\par\medskip\Needspace{5\baselineskip}\noindent\textbf{解}\quad\ignorespaces}{\par\medskip}
\newenvironment{drillsolution}[1]{\subsubsection*{练习\ref{#1}详解}\phantomsection\label{sol:#1}\ignorespaces}{\par\medskip}
\newcommand{\dd}{\mathop{}\!\mathrm d}
\tikzset{axis/.style={->,>=Stealth,line width=.45pt},world/.style={line width=.8pt},flow/.style={->,>=Stealth,line width=.7pt},ray/.style={line width=.7pt,->,>=Stealth},guide/.style={dashed,line width=.35pt},dim/.style={<->,>=Stealth,line width=.35pt},every node/.style={font=\small},force/.style={->,>=Stealth,very thick},gas/.style={fill=blue!7}}
\newsavebox{\LectureFigureBox}\newcommand{\fig}[3]{\begin{figure}[H]\centering\sbox{\LectureFigureBox}{#3}\ifdim\wd\LectureFigureBox>\linewidth\resizebox{\linewidth}{!}{\usebox{\LectureFigureBox}}\else\usebox{\LectureFigureBox}\fi\caption{#2}\label{#1}\end{figure}}
\setlist[enumerate]{leftmargin=2em,itemsep=2pt,topsep=3pt}
\fancyhead[RE]{\small 第1讲\quad 物质的热性质}
% Common visual style for all first-round physics lecture projects.
% Each project includes an identical local copy of this file.
\geometry{paperwidth=176mm,paperheight=250mm,inner=21mm,outer=18mm,top=20mm,bottom=20mm,headheight=14pt,headsep=6mm,footskip=10mm}
\hypersetup{hidelinks}
\setlength{\parindent}{2em}\setlength{\parskip}{3pt}\setlength{\emergencystretch}{2em}
\linespread{1.18}\raggedbottom\setcounter{tocdepth}{1}
\widowpenalty=10000\clubpenalty=10000\displaywidowpenalty=10000
\setlength{\intextsep}{8pt plus1pt minus1pt}
\makeatletter
\g@addto@macro\normalsize{\setlength{\abovedisplayskip}{8pt plus2pt minus2pt}\setlength{\belowdisplayskip}{8pt plus2pt minus2pt}\setlength{\abovedisplayshortskip}{0pt plus2pt}\setlength{\belowdisplayshortskip}{6pt plus2pt minus2pt}}
\makeatother
\providecommand{\Needspace}[1]{}
\newdimen\LectureAvailableSpace
\renewcommand{\Needspace}[1]{\par\penalty-100\begingroup\LectureAvailableSpace=\pagegoal\advance\LectureAvailableSpace by-\pagetotal\ifdim\LectureAvailableSpace<#1\relax\ifdim\pagetotal>0pt\newpage\fi\fi\endgroup}
\renewenvironment{solution}{\par\smallskip\Needspace{3\baselineskip}\noindent\textbf{解}\quad\ignorespaces}{\par\medskip}
\ctexset{chapter={name={第,讲},number=\chinese{chapter},format=\LARGE\bfseries\centering,beforeskip=16pt,afterskip=22pt},section={format=\Large\bfseries},subsection={format=\large\bfseries},subsubsection={format=\normalsize\bfseries,beforeskip=10pt,afterskip=4pt}}
\setlist[enumerate]{leftmargin=2em,itemsep=2pt,topsep=3pt}
\pagestyle{fancy}\fancyhf{}\fancyhead[LE,RO]{\small\thepage}\fancyhead[LO]{\small 物理竞赛一轮讲义}
\renewcommand{\headrulewidth}{0pt}
\fancypagestyle{plain}{\fancyhf{}\fancyfoot[C]{\small\thepage}\renewcommand{\headrulewidth}{0pt}}
\renewenvironment{drillsolution}[1]{\par\Needspace{4\baselineskip}\subsubsection*{练习\ref{#1}详解}\phantomsection\label{sol:#1}\nobreak\noindent\textbf{解}\quad\ignorespaces}{\par\medskip}
\ifdefined\exercisesolution
\renewenvironment{exercisesolution}[1]{\par\Needspace{4\baselineskip}\subsubsection*{习题\ref{#1}详解}\phantomsection\label{sol:#1}\nobreak\noindent\textbf{解}\quad\ignorespaces}{\par\medskip}
\fi
\newcommand{\LectureAnswerSection}{\par\section*{练习解答}\addcontentsline{toc}{section}{练习解答}}
\newcommand{\LectureHeader}[2]{\fancyhead[RE]{\small #1\quad #2}}
\newcommand{\LectureCover}[3]{\begin{titlepage}\centering\vspace*{32mm}{\Huge\bfseries #1\par}\vspace{12mm}{\Large 物理竞赛一轮讲义\par}\vspace{18mm}{\LARGE #2\\[5mm]#3\par}\vfill{\large 2026\par}\end{titlepage}}

\LectureHeader{第一讲}{物质的热性质}
\begin{document}\frontmatter\LectureCover{热学}{第一讲}{物质的热性质}\tableofcontents\mainmatter\setcounter{chapter}{0}\chapter{物质的热性质}
\section{热膨胀与温度测量}
\subsection{从温度变化到长度变化}
温度是描述物体热状态的量。两个物体接触后不再发生净热传递，称为达到热平衡；热平衡时温度相同。物体含有多少内能，还与物质的种类和多少有关，不能用内能大小定义温度。热量则是在温差作用下传递的能量，不是物体“含有”的状态量。

摄氏温度与热力学温度的换算写为
\[\frac T{\mathrm K}=\frac t{{}^\circ\mathrm C}+273.15.\]
温差的数值以K或摄氏度差表示相同，但温度之比必须用K。例如气态方程中的 $T$ 不能替换为摄氏温度；在热膨胀公式中，若只出现 $\Delta T$，则两种温差单位的数值可直接对应。部分题为计算方便取 $273$，应遵从该题约定，不能在同一计算中混用。


温度升高时，固体微观粒子的平均间距通常增大，宏观长度随之改变。在一段温区内，实验发现长度变化近似正比于原长与温差：

\[\Delta L=\alpha L_0\Delta T,\qquad L=L_0(1+\alpha\Delta T).\]
其中 $L_0$ 是基准温度 $T_0$ 下的长度，$\alpha$ 的单位是 $\mathrm{K^{-1}}$。此公式含两个近似：材料可自由膨胀，且该温区内 $\alpha$ 变化不明显。若杆被刚性支架固定，长度未必按上式实际增加，而可能产生热应力，此时还须结合弹性关系求解。

\fig{fv:h1:expansion}{自由杆的原长与伸长分开标记；图中伸长夸张绘制}{\begin{tikzpicture}[x=1cm,y=.9cm]
\draw[thick](0,1)rectangle(4,1.4);\draw[thick](0,0)rectangle(5,.4);\draw[dashed](4,-.3)--(4,1.7);
\draw[<->](0,2)--(4,2)node[midway,above]{$L_0$};\draw[<->](4,-.6)--(5,-.6)node[midway,below]{$\Delta L$};
\node[right]at(5.3,1.2){$T_0$};\node[right]at(5.3,.2){$T$};\end{tikzpicture}}
图\ref{fv:h1:expansion}中两次长度共用同一左端基准，虚线保留原右端的位置，向右多出的一段才是$\Delta L$。相对伸长应除以原长$L_0$；不能将升温后的$L$误作基准长度。若两端都受固定限制，图示自由伸长便不能实际发生，应另列约束和弹性关系。


从相对变化看，$\Delta L/L_0=\alpha\Delta T$。多数金属的 $\alpha$ 很小，所以在常见温差下可以保留一阶小量。若长度的变化进入周期、密度等非线性表达式，应先写精确的量之间关系，再统一展开，避免对不同量随意作同一种近似。

热膨胀系数给出的是材料的响应，题目要求的往往是装置的响应。以摆钟为例，温度先改变摆长，摆长再改变一次摆动所需的时间，最后才改变一天内完成的摆动次数。这三个环节应分别列式。若把“长度增加万分之几”直接当成“钟少走万分之几”，便忽略了周期对长度的平方根依赖。

对于随长度的幂次变化的量$Y\propto L^q$，令$\varepsilon=\alpha\Delta T$，先写$Y/Y_0=(1+\varepsilon)^q$。对$F(\varepsilon)=(1+\varepsilon)^q$求导，
\[
F'(\varepsilon)=q(1+\varepsilon)^{q-1},\qquad
F(0)=1,\quad F'(0)=q.
\]
在$|\varepsilon|\ll1$时，一阶展开$F(\varepsilon)=F(0)+F'(0)\varepsilon+O(\varepsilon^2)$才给出$\Delta Y/Y_0\simeq q\alpha\Delta T$。比较实际时间与仪器读数时，还须由计数方式判断应用的是周期本身，还是周期的倒数。

若膨胀系数显著依赖温度，可用微分定义 $\alpha=(1/L)(\dd L/\dd T)$。将变量分离为$\dd L/L=\alpha(T)\dd T$，并把基准状态$L(T_0)=L_0$作为积分下限，得到
\[
\int_{L_0}^{L}\frac{\dd\ell}{\ell}
=\int_{T_0}^{T}\alpha(\vartheta)\,\dd\vartheta,
\qquad
[\ln\ell]_{L_0}^{L}=\ln\frac{L}{L_0}.
\]
对两边取指数，才有$L/L_0=\exp[\int_{T_0}^{T}\alpha(\vartheta)\dd\vartheta]$。常系数时右侧积分为$\alpha(T-T_0)$，再展开$e^{\alpha\Delta T}=1+\alpha\Delta T+O((\alpha\Delta T)^2)$，便得到本讲使用的一阶模型。




\subsection{面积、体积、孔洞与复合结构}
对各向同性材料，三个方向按同一长度比例变化。因此面积与体积分别满足
\[\frac A{A_0}=(1+\alpha\Delta T)^2\simeq1+2\alpha\Delta T,\qquad
\frac V{V_0}=(1+\alpha\Delta T)^3\simeq1+3\alpha\Delta T.\]
令$\varepsilon=\alpha\Delta T$，平方与立方的实际展开为$1+2\varepsilon+\varepsilon^2$和$1+3\varepsilon+3\varepsilon^2+\varepsilon^3$。舍去二阶及更高项，才得到面积膨胀系数约为$2\alpha$、体膨胀系数约为$3\alpha$；前提是材料各向同性。质量不变时，先用$\rho V=\rho_0V_0$得$\rho/\rho_0=(1+\varepsilon)^{-3}$，再在上面的幂函数展开中取$q=-3$，得到$\rho\simeq\rho_0(1-3\alpha\Delta T)$。

孔洞随整体的几何尺度扩大。可把孔暂时用同种材料填满，整体自由、均匀膨胀后再取出填充部分；孔的边界会移动到相同的放大位置。这个论证要求自由膨胀，不能直接用于外边界被夹紧的金属环。

也可以从坐标来理解孔洞膨胀。若以任意一点为原点，自由均匀膨胀使每个材料点的坐标都乘 $1+\alpha\Delta T$。孔边界上的点服从相同的位移规律，因此孔半径也乘这一因子。它说明“孔变大”是整个位移场的结果，与孔内有没有填充物无关。若外缘受固定约束，位移场已经改变，便需要重新求解，不能仍沿用自由膨胀的结论。

\fig{fig:h1:holeconstraint}{自由圆环的内外边界按同一比例放大；两端固定的杆则建立压应力}{\begin{tikzpicture}

\draw[guide](0,0)circle(.8) (0,0)circle(1.6);\draw[blue,thick](0,0)circle(.92) (0,0)circle(1.84);\fill(0,0)circle(.04)node[below left]{$O$};\draw[flow](1.6,0)--(1.84,0);\draw[flow](0,.8)--(0,.92);\node at(0,-2.4){$a'=(1+\alpha\Delta T)a$};
\begin{scope}[xshift=7cm]\draw[world](0,-.8)--(0,.8) (4,-.8)--(4,.8);\draw[world](0,-.2)rectangle(4,.2);\draw[force](0,0)--(1,0)node[above]{$F$};\draw[force](4,0)--(3,0)node[above]{$F$};\node at(2,1.2){两端固定};\node at(2,-1.4){$\varepsilon_{\rm total}=0$};\end{scope}

\end{tikzpicture}}
图\ref{fig:h1:holeconstraint}左图中，材料点的位移沿半径向外，内边界并没有向孔内挤入。右图的杆却不能伸长：若再采用线弹性模型、杨氏模量为$E$，总应变由热应变与弹性应变相加，

\[
 0=\alpha\Delta T+\frac{\sigma_m}{E},\qquad
 \sigma_m=-E\alpha\Delta T.
\]
这里约定拉应力为正，故升温后的应力为压应力。这个结果来自固定端的几何限制；自由杆只有热伸长，不能同时再给它施加上述固定端应力。



复合杆、双金属结构和有间隙的装置应逐段写变化量，再列几何约束。串联自由杆的总长变化为 $\sum_i\alpha_iL_i\Delta T$；两个端点的间距变化则可能是两段伸长之差。解题时应在图上标出比较的两个端点，不能看到“加热”就把所有变化量相加。

若尺本身膨胀，每个刻度对应的实际长度也随之增加，读数应由“被测长度除以实际刻度长度”求出。尺与被测物同时升温时，一阶读数变化取决于两者膨胀系数之差，而不是只取其中一个系数。例\ref{t1-1}和练习\ref{t1-10}分别讨论时间测量与孔洞几何。

\section{表面张力与弯曲液面}
\subsection{表面张力的方向与两个表面的计数}
液体表层分子的环境不同于内部，形成液面需要能量。在给定温度和界面成分下，液面有减小面积的趋势。沿液面内任取一条边界线，邻侧液面对该线的作用力位于液面切平面内，方向垂直于该边界线。

对长度为 $\ell$ 的直边，一个液面贡献 $\sigma\ell$ 的拉力。肥皂膜由两面构成，若两面表面张力系数均为 $\sigma$，活动边受到的总拉力为 $2\sigma\ell$。这里的“两个表面”不能误解成边界线的左右各算一次；应直接数同一薄膜的正反两个气液界面。

\fig{fig:film-work}{双面膜活动边的受力与面积变化}{\begin{tikzpicture}[x=1cm,y=1cm]
\fill[blue!8] (0,0) rectangle (3,2);\draw[world] (0,2)--(0,0)--(3,0) (0,2)--(3,2);\draw[very thick] (3,0)--(3,2);\draw[guide] (4,0)--(4,2);\draw[force] (3,1)--(1.8,1) node[midway,above]{$2\sigma\ell$};\draw[force] (3,.45)--(4.3,.45)node[right]{$F_{\text{外}}$};\draw[flow] (3,2.5)--(4,2.5) node[midway,above]{$\dd x$};\draw[dim] (4.6,0)--(4.6,2) node[midway,right]{$\ell$};\begin{scope}[xshift=6cm]\draw[blue,thick](0,.5)--(1.5,.5) (0,.3)--(1.5,.3);\node at(.75,1.3){膜厚夸张的截面};\node[right]at(1.5,.8){正面};\node[right]at(1.5,0){反面};\end{scope}
\end{tikzpicture}}

图\ref{fig:film-work}的活动边向右移动 $\dd x$ 时，每面面积增加 $\ell\dd x$，总面积增加 $2\ell\dd x$。外力准静态做功 $2\sigma\ell\dd x$，恰等于表面自由能的增加 $\sigma\dd A_{\text{总}}$。因此 $\sigma$ 既可用 $\mathrm{N/m}$ 表示，也可用 $\mathrm{J/m^2}$ 表示；面积法与受力法是相互校验的两种方法。

恒温且 $\sigma$ 近似不变时，$\Delta F_{\text{表}}=\sigma\Delta A_{\text{总}}$。若温度变化明显或界面成分变化，$\sigma$ 也会变化，不能直接把全部表面能变化写成同一常数乘面积差。

\subsection{由局部受力确定曲率与压差}
受均匀膜张力作用的轻细线，平衡时曲率与张力有关。在光滑圆弧上取长度 $\dd s=r\dd\varphi$ 的微元，线两端张力的法向合力为 $\tau\dd\varphi$；双面膜的横向拉力为 $2\sigma\dd s$。因此

\[\frac\tau r=2\sigma.\]
\fig{fig:h1:lineelement}{弧形细线微元的两端张力合力向曲率中心，膜力沿法向向外}{\begin{tikzpicture}

\coordinate(A)at(1.691447,-.615636);\coordinate(B)at(1.691447,.615636);\draw[blue,very thick](A)arc(-20:20:1.8);\draw[guide](0,0)--(A) (0,0)--(B);\fill(0,0)circle(.04)node[left]{$O$};\node[below]at(.9,-.33){$r$};\draw(.6,-.218382)arc(-20:20:.638507);\node at(.95,0){$\dd\varphi$};
\draw[force](A)--(1.281023,-1.743267)node[below]{$\tau$};\draw[force](B)--(1.281023,1.743267)node[above]{$\tau$};\draw[force](1.8,0)--(3.2,0)node[right]{$2\sigma\dd s$};\node[right]at(1.9,.65){$\dd s=r\dd\varphi$};

\end{tikzpicture}}
图\ref{fig:h1:lineelement}把“沿细线的张力$\tau$”与“膜对细线的分布拉力$2\sigma$”分开。前者在两端分别沿切线，后者垂直线元；两端张力的方向差才提供$\tau\dd\varphi$，不能把两个张力大小直接相加为$2\tau$。



上述微元平衡还说明自由段为什么是圆弧。轻线处处无切向分布力，张力的大小沿线不变；膜的单位长度拉力也是常量，所以 $1/r=2\sigma/\tau$ 处处相同。曲率恒定的平面曲线就是圆弧。应先确定这条局部关系，再用全局长度和容许区域确定圆弧的位置。若所需圆弧越过容器边界，继续增加的线长只能通过接触段等方式容纳。

局部关系只决定弯曲段的半径，还必须满足端点、长度和边界接触条件。例\ref{t1-2}中长线不能全部画成自由优弧，应检查上下边界是否阻止圆弧继续延伸。

球形液面也可由整体受力求附加压强。图\ref{fig:laplace}取上半球，向上的压差合力作用在圆形截面，沿边缘向下的表面张力作用于整个圆周。取半球，截面上的压差合力是 $\Delta p\,\pi r^2$，圆周表面张力合力是 $2\pi r\sigma$，所以

\[\boxed{\Delta p=\frac{2\sigma}{r}}.\]
对液滴或液体中的单界面气泡，内侧压强均比紧邻外侧高 $2\sigma/r$；薄肥皂泡有两个表面，近似为 $\Delta p=4\sigma/r$。判断压差符号时，应先画出表面张力对半球的作用方向，再列平衡方程。不能单凭“里面是气体还是液体”决定压差正负。

\fig{fig:laplace}{半球截面上的压差合力，与截面圆周上的表面张力相平衡}{\begin{tikzpicture}

\fill[blue!8](0,0)arc(180:0:1.2)--cycle;\draw[world](0,0)arc(180:0:1.2);\draw[guide](0,0)--(2.4,0);\draw[force](1.2,.2)--(1.2,2)node[above]{$(p_{\rm in}-p_{\rm out})\pi r^2$};\draw[force](0,0)--(0,-.9);\draw[force](2.4,0)--(2.4,-.9);\node at(1.2,-1.35){圆周总拉力$2\pi r\sigma$};\draw[dim](1.2,-.4)--(2.4,-.4)node[midway,below]{$r$};

\end{tikzpicture}}

\subsection{润湿、接触角与毛细高度}
液体与固体接触时，界面形状取决于几种界面的相互作用。接触角 $\theta$ 在液体一侧量取：$\theta<90^\circ$ 为润湿，弯月面通常凹向气体；$\theta>90^\circ$ 为不润湿，弯月面通常凸向气体。画接触角时，必须使用液面的切线，不能用连接两端的水平弦代替切线。

取半径 $r$ 的润湿毛细管中高于槽面的液柱，如图\ref{fig:h1:capillarybalance}左侧；右侧将这段液柱的剩余竖直作用单独画出。接触线总长 $2\pi r$，表面张力的竖直合力为 $2\pi r\sigma\cos\theta$。忽略弯月面体积，液柱重力为 $\rho g\pi r^2h$，故

\[\boxed{h=\frac{2\sigma\cos\theta}{\rho gr}}.\]
\fig{fig:h1:capillarybalance}{润湿液柱的高度与等效竖直受力；左图以$\theta=60^\circ$示意}{\begin{tikzpicture}

\draw[world](0,-.5)--(0,2.8) (.8,-.5)--(.8,2.8);\draw[guide](-.8,0)--(2,0)node[right]{槽面};\fill[blue!8](0,0)--(0,1.7)arc(240:300:.8)--(.8,0)--cycle;\draw[blue,thick](0,1.7)arc(240:300:.8);
\draw[guide](0,1.7)--(.6,1.35359);\draw(0,1.4)arc(-90:-30:.3);\node[right]at(.18,1.2){$\theta$};\draw[dim](1.5,0)--(1.5,1.6)node[midway,right]{$h$};\draw[force](0,1.7)--(-.8,2.16188);\draw[force](.8,1.7)--(1.6,2.16188);
\begin{scope}[xshift=5.5cm,yshift=.8cm]\fill(0,0)circle(.07);\draw[force](0,0)--(0,1.4)node[above]{$2\pi r\sigma\cos\theta$};\draw[force](0,0)--(0,-1.4)node[below]{$\rho g\pi r^2h$};\end{scope}

\end{tikzpicture}}
图\ref{fig:h1:capillarybalance}左图接触线处的箭头沿界面切线，两个箭头是整个圆周拉力的截面示意。以方位角$\varphi$参数化接触圆周，$\dd\ell=r\dd\varphi$；每段竖直拉力为$\sigma\cos\theta\,\dd\ell$，所以
\[
F_z=\int_0^{2\pi}\sigma r\cos\theta\,\dd\varphi
=\sigma r\cos\theta[\varphi]_0^{2\pi}
=2\pi r\sigma\cos\theta.
\]
这才是右图所标的总力。忽略弯月面体积时，液柱体积近似为$\pi r^2h$，所以周长与面积不同的尺度依赖给出$h\propto1/r$。



$h$ 是带符号的液面相对高度，润湿时为正，不润湿时为负。由公式可见管越细、毛细效应越显著；这来自张力随周长变化、重力随面积变化，而不是细管“吸力更强”的额外机制。

毛细高度也可由压强验证。润湿液面的曲率半径约为 $r/\cos\theta$，气体一侧压强比液面下方高 $2\sigma\cos\theta/r$。从槽面沿液体向上走 $h$，液体压强降低 $\rho gh$，两种压差相等，仍得到同一高度。受力法取的是整段液柱，压强法比较的是液面两侧；把两种方法的力或压差重复相加，会把同一表面作用算两次。

上述推导假设槽面近似不动、管径远小于毛细高度，且液面已达到平衡。若两管连接有限容积液槽，应再列液体体积守恒。若存在快速流动，黏性阻力与动态接触角也会影响过程，静态毛细高度不能代替流量公式。

植物输水中，毛细作用可以说明局部弯月面压差，但完整运输还涉及蒸腾引起的水势差及连续水柱等条件。将静态毛细高度直接当作植物输水的全部解释，会混淆局部平衡与持续输运。

\fig{fig:contact-angle-sign}{接触角从液体一侧量取，界面切线决定张力竖直分量；弯月面作局部示意}{
\begin{tikzpicture}[x=1.1cm,y=1.1cm]
\foreach \s in {0,4.4}{\begin{scope}[xshift=\s cm]
\draw[world] (0,-1.5)--(0,1.5) (.8,-1.5)--(.8,1.5);
\draw[guide] (-.8,0)--(1.7,0);\node[right] at (1.7,0){槽面};
\end{scope}}
\fill[blue!10] (0,-1.5)--(0,.8) .. controls (.2,.6) and (.6,.6) .. (.8,.8)--(.8,-1.5)--cycle;
\draw[world] (0,.8) .. controls (.2,.6) and (.6,.6) .. (.8,.8);
\draw[guide] (0,.8)--(.48,.32);\draw[force] (0,.8)--(-.55,1.35);\draw[force] (.8,.8)--(1.35,1.35);
\draw (0,.56) arc(-90:-45:.24);\node at (.23,.48){$\theta$};
\node[align=center] at (.4,-2){$\theta<90^\circ$\\$\cos\theta>0,\ h>0$};
\begin{scope}[xshift=4.4cm]
\fill[blue!10] (0,-1.5)--(0,-.8) .. controls (.2,-.6) and (.6,-.6) .. (.8,-.8)--(.8,-1.5)--cycle;
\draw[world] (0,-.8) .. controls (.2,-.6) and (.6,-.6) .. (.8,-.8);
\draw[guide] (0,-.8)--(.48,-.32);\draw[force] (0,-.8)--(-.5,-1.3);\draw[force] (.8,-.8)--(1.3,-1.3);
\draw (0,-1.04) arc(-90:45:.24);\node at (.32,-.99){$\theta$};
\node[align=center] at (.4,-2){$\theta>90^\circ$\\$\cos\theta<0,\ h<0$};
\end{scope}
\end{tikzpicture}}
图\ref{fig:contact-angle-sign}中，箭头表示接触线处作用于液体的张力方向，沿气液界面的切线指向边界。润湿时其竖直分量向上，不润湿时向下；改变的不是 $\sigma$ 的正负，而是 $\cos\theta$ 的正负。对不润湿的下降液面，宜用压强法统一处理：
\[
 p_0-p_{\text{液面下}}=\frac{2\sigma\cos\theta}{r},
 \qquad p_0-p_{\text{液面下}}=\rho gh.
\]
两式都按同一方向取压差，因而给出同一个带符号高度公式。此时不能把“槽面以上的一段液柱”当作实际存在的正高度液柱，再给其重力随意加负号。图中的槽面是固定参考高度，$h<0$ 表示管内液面低于它。

\section{理想气体与液柱压强}
\subsection{状态量、绝对压强与理想气体近似}
描述定量平衡气体通常使用压强 $p$、体积 $V$ 和绝对温度 $T$。压强是气体对边界单位面积的法向作用，体积是气体实际占据的空间，温度反映热运动状态。理想气体模型忽略分子自身占有的体积及分子间势能；气体稀薄且远离凝结时，通常较为适用。

理想气体状态方程为
\begin{equation}\label{eq:t1state}
pV=nRT=NkT,\qquad n=\frac m\mu,\quad N=nN_A,\quad k=\frac R{N_A}.
\end{equation}
其中 $n$ 是物质的量，$N$ 是分子数，$\mu$ 是摩尔质量，三者不能混写。使用SI单位时 $p$ 用Pa，$V$ 用 $\mathrm{m^3}$，$\mu$ 用kg/mol，$R=8.314\,\mathrm{J/(mol\,K)}$。

式\eqref{eq:t1state}中的 $p$ 是相对于真空的绝对压强。若压力计给的是表压 $p_g$，应换成 $p=p_0+p_g$；只有在求活塞净受力时，才可直接用两侧压差乘面积。常见的错误是把表压代入气态方程，却在另一状态使用绝对压强。


固定气体量时，状态方程给出三种基本过程：等温 $pV=\text{常量}$，等容 $p/T=\text{常量}$，等压 $V/T=\text{常量}$。每一种关系都还包含“气体量固定”这一条件，漏气、充气、凝结或化学反应时必须重新检查。

在 $p$--$V$ 图上，等温线是双曲线，等容线竖直、等压线水平。在 $p$--$T$ 图上，等容线是过原点的射线，斜率 $nR/V$；在 $V$--$T$ 图上，等压线是过原点的射线，斜率 $nR/p$。若横轴采用摄氏温度，直线一般不再过坐标原点。

\fig{fig:h1:stateplots}{固定气体量时的过程图；温度坐标必须使用热力学温度}{\begin{tikzpicture}

\draw[axis](0,0)--(3.7,0)node[right]{$V$};\draw[axis](0,0)--(0,3.2)node[above]{$p$};\draw[blue,thick,domain=.8:3.3,samples=80]plot(\x,{2/\x});\draw[red,thick](.6,1)--(3.3,1);\draw[green!50!black,thick](2,.4)--(2,2.8);\fill(2,1)circle(.05);\node[blue]at(3,2.2){等温};\node[red]at(3,1.35){等压};\node[green!50!black]at(1.8,3.2){等容};
\begin{scope}[xshift=5.3cm]\draw[axis](0,0)--(3.6,0)node[right]{$T$};\draw[axis](0,0)--(0,3.2)node[above]{$p$};\draw[blue,thick](0,0)--(3.1,2.3);\node at(1.8,3.1){$V$固定};\node at(1.7,-1){斜率$nR/V$};\end{scope}
\begin{scope}[xshift=10.5cm]\draw[axis](0,0)--(3.6,0)node[right]{$T$};\draw[axis](0,0)--(0,3.2)node[above]{$V$};\draw[red,thick](0,0)--(3.1,2.3);\node at(1.8,3.1){$p$固定};\node at(1.7,-1){斜率$nR/p$};\end{scope}

\end{tikzpicture}}
图\ref{fig:h1:stateplots}左图三种过程可经过同一状态，但它们的后续路径不同。中图、右图的直线通过$T=0$的坐标原点是理想气体关系的数学外推；实际气体在极低温可能凝结，不能据外推断定模型在那里仍适用。用摄氏温度作横坐标时，只平移横轴，不能仍强令直线经过新原点。



两个平衡态之间总可写 $p_1V_1/T_1=p_2V_2/T_2$，但这只是端点关系，不能据此断定中间全过程等压、等温或等容。例\ref{t1-5}中三个阶段共用同一端点状态方程，却分别由不同机械约束确定路径。

\subsection{液柱压强与多液面联动}
静止液体中，取底面积$S$、高度$\dd z$的竖直薄层，向上为正。底面压力、顶面压力与重力平衡给出$p(z)S-p(z+\dd z)S-\rho gS\dd z=0$，除以$S\dd z$并取极限得$\dd p/\dd z=-\rho g$。若密度均匀，从$z_1$积分到$z_2$，
\[
\int_{p_1}^{p_2}\dd p=-\rho g\int_{z_1}^{z_2}\dd z,
\qquad p_2-p_1=-\rho g(z_2-z_1),
\]
故$p+\rho gz=\text{常量}$。沿同一连续液体向下走高度$h$，压强增加$\rho gh$；向上走则减小。同一水平面等压的结论只能用于同种、连通且静止的液体，不能跨过隔板或气体直接使用。

\fig{fig:hydrostatic}{同一液体内从开口端到气体端的压强传递}{\begin{tikzpicture}[x=1cm,y=1cm]
\draw[world] (0,3)--(0,0)--(3,0)--(3,2.5) (.4,3)--(.4,.4)--(2.6,.4)--(2.6,2.5);\draw[thick] (0,2.3)--(.4,2.3) (2.6,1.2)--(3,1.2);\draw[guide] (-.3,1.2)--(3.3,1.2);\node[above] at (.2,3){$p_0$};\node[above] at (2.8,2.5){$p$};\draw[dim] (3.8,1.2)--(3.8,2.3) node[midway,right]{$\Delta h$};\node at (1.6,-.7){$p=p_0+\rho g\Delta h$};
\end{tikzpicture}}

图\ref{fig:hydrostatic}中先从左液面向下到辅助水平面，再沿该水平面走到右液面，得到 $p=p_0+\rho g\Delta h$。若气体侧液面更高，$\Delta h$ 应取负。建议始终保留高度坐标，最后再代入数值，避免凭记忆写出错误的正负号。

液面发生位移时，还要满足液体体积守恒。若一U管两臂截面为 $S_1,S_2$，带符号液面位移满足 $S_1\delta h_1+S_2\delta h_2=0$。等截面才有两侧位移大小相等；由此得到的液面高度差变化常为单侧位移的两倍。例\ref{t1-3}用两组体积守恒再联系同一封闭气体，例\ref{t1-4}还须把管顶固定转化为空气柱长度变化。


理想混合气体中，每种组分可用其物质的量定义分压：$p_iV=n_iRT$，总压 $p=\sum_i p_i$。分压对应同一最终体积与温度，不能把不同容器、不同温度的初压直接相加。

若多个刚性容器等温连通且连通管体积可忽略，气体总量守恒，得到 $p_f\sum V_i=\sum p_iV_i$。若初温不同，应先列总内能守恒或与环境的能量交换，求出末温后再用状态方程。饱和蒸气还可能凝结或汽化，气相物质的量未必固定，这将在第四讲单独处理。

\section{活塞、轻绳与分段过程}
\subsection{机械平衡怎样确定气体路径}
状态方程只能联系 $p,V,T$，还需要机械条件确定过程。轻活塞在恒定外压下自由、缓慢运动时，两侧压强相等，因而气体可近似等压；重活塞竖直放置时，平衡压强为 $p=p_0+mg/S$；水平弹簧活塞则满足 $p=p_0+kx/S$，通常不是等压过程。

\fig{fig:piston-forces}{负载与挡板反力共同确定活塞压强；挡板反力只在接触时存在}{\begin{tikzpicture}

\foreach \j/\lab in {0/轻活塞,1/重活塞,2/重活塞接触上挡板}{\begin{scope}[xshift={4.6*\j*1cm}]
\draw[world](0,2.7)--(0,0)--(1.8,0)--(1.8,2.7);\fill[blue!8](0,0)rectangle(1.8,1.2);\draw[very thick](0,1.2)--(1.8,1.2);\node at(.9,.5){$p$};\node at(.9,2.4){$p_0$};\node at(.9,-.55){\lab};
\draw[force](.4,1.2)--(.4,2.1)node[left]{$pS$};\draw[force](1.4,2.1)--(1.4,1.2);\end{scope}}
\node at(1.8,1.8){$p_0S$};\node[align=left]at(6.9,1.7){$p_0S$\\$+mg$};\node[align=left]at(11.5,1.7){$p_0S$\\$+mg+N$};\draw[world](9.2,1.2)rectangle(11,1.35);\draw[world,very thick](9.2,1.35)--(11,1.35);
\node at(.9,-1.25){$p=p_0$};\node at(5.5,-1.25){$p=p_0+mg/S$};\node at(10.1,-1.25){$p=p_0+(mg+N)/S$};

\end{tikzpicture}}

图\ref{fig:piston-forces}显示同一气体方程对应不同机械路径。上挡板接触时，挡板可向下提供支持力，因此气体压强可超过自由平衡值而活塞仍不动。挡板撤去或支持力降为零，才切换到自由运动段。


轻绳只能拉，不能推；张力 $\tau\ge0$。计算得到负张力，通常表示“绳绷直”的假设失效，应改用绳松弛的约束。接触支持力同样只能指向阻止物体穿过边界的方向。若某段计算需要挡板把活塞拉回来，该段就不成立。

分段分析可按照以下顺序进行：先假定当前接触和绳状态；对每个活塞列受力平衡；用位置写气体体积；结合气态方程求状态变化；最后检查接触、张力和几何边界。某个不等式首次达到等号，便是下一段的起点。

在同一段中应消去不必要的未知量。例如两个活塞由绳连接且同时自由，分别列平衡后消去张力，往往可以先确定压强；再由绳长守恒把两个位移归结为一个坐标。台阶管中，两个活塞同向移动也可能使气体体积减小，体积应按各段截面积乘长度重新求和。

“接触挡板”与“自由运动”应由受力连续地衔接。例如重活塞靠在上挡板上时，$N=(p-p_0)S-mg\ge0$，压强可以高于自由平衡值。冷却使 $N$ 逐渐减小，只有减到零，活塞才开始离开；若继续把它当作固定边界，就会算出负的 $N$。反过来，自由运动段中位置达到挡板是几何临界条件，此刻支持力尚为零，随后升温才逐渐建立支持力。临界状态同时满足相邻两段的条件，因此可用它检查分段公式是否连续。

同理，对连接两活塞的绳，应先标清气体在哪一侧。气体对两活塞的压力未必同向，两个压差力也未必使用同一面积。先逐个活塞列式再消去张力，才能判断共同移动是等压过程还是变压过程。

例\ref{t1-8}训练“自由平移—接触台阶”的转换；例\ref{t1-9}包含“共同移动—固定体积—绳松弛”三段。解答不能只给一条折线，还要说明每个转折点对应装置中的哪个事件，临界张力与支持力是否符合方向约束。

\subsection{可逆绝热过程中的临界状态}
例\ref{t1-7}还需把受力确定的临界压强换算成温度。这里先使用定量理想气体的可逆绝热关系；其完整推导与自由膨胀的比较见热学第三讲。若热容恒定、过程缓慢无摩擦且无热交换，初态为 $(p_1,T_1)$、临界态为 $(p_2,T_2)$，则

\[
 \frac{T_2}{T_1}=\left(\frac{p_2}{p_1}\right)^{(\gamma-1)/\gamma},
 \qquad \gamma=\frac{C_p}{C_V}>1.
\]
先对气缸壳体列竖直受力平衡，令地面支持力首次降为零，求得 $p_2$；再将这个压强代入温度比。这样，离地条件由机械约束确定，达到该状态所需的温度由热学路径确定，两步各有独立作用。压缩时 $p_2>p_1$会给出 $T_2>T_1$，也可用来检查指数的正负。仅有“缸壁绝热”而过程不满足上述条件时，不能直接使用这个幂律。


\section{例题与分析}
\par\medskip\noindent\begin{minipage}{\textwidth}
\begin{example}[钢摆钟的日差]\label{t1-1}
一钢制单摆在 $0^\circ\mathrm C$ 时周期 $P_0=1.000\,\mathrm s$，线膨胀系数 $\alpha=1.2\times10^{-5}\,\mathrm{K^{-1}}$。温度升至 $25^\circ\mathrm C$，重力加速度不变。求新周期，以及实际经过一天时钟面少走的时间。见图\ref{fig:t1-1q}。
\end{example}
\fig{fig:t1-1q}{钢摆钟的日差：题设}{\begin{tikzpicture}[x=1cm,y=1cm]
\draw[world] (0,2)--(1,2);\draw[world] (.5,2)--(.5,-.4);\fill (.5,-.4) circle(.15);\draw[dim] (1.2,2)--(1.2,-.4) node[midway,right]{$L_0$};\draw[world] (4,2)--(5,2);\draw[world] (4.5,2)--(4.5,-.6);\fill (4.5,-.6) circle(.15);\draw[dim] (5.2,2)--(5.2,-.6) node[midway,right]{$L$};\node at (.5,-1){$0^\circ\mathrm C$};\node at (4.5,-1){$25^\circ\mathrm C$};
\end{tikzpicture}}
\end{minipage}
\begin{solution}\textbf{先求自由热伸长，再由摆长求周期}\par\nobreak
图\ref{fig:t1-1q}中两根摆线表示同一钢摆在不同温度下的长度。支点不变、重力加速度不变，温度只通过摆长影响周期。由自由线膨胀关系，
\[L=L_0(1+\alpha\Delta T),\qquad \alpha\Delta T=1.2\times10^{-5}\times25=3.0\times10^{-4}.\]
将它代入单摆的小振幅周期公式，而不是直接把长度变化率作为周期变化率，得
\[\frac P{P_0}=\sqrt{\frac L{L_0}}=\sqrt{1.00030},\qquad \boxed{P\simeq1.000150\,\mathrm s}.\]


\par\medskip\Needspace{4\baselineskip}\noindent\textbf{用振动计数定义钟面时间}\par\nobreak
接下来需要区分实际经过的时间与钟面计数。设实际经过 $t$，摆共完成 $N=t/P$ 次周期，钟的刻度机构仍按每次 $P_0$ 秒累计，因此钟面显示 $t_{\text{钟}}=NP_0$。图\ref{fig:t1-1s}用相邻峰的间隔表示周期，并用同一条竖线比较实际经过相同时间后的振动计数；周期增长后计数较少。\fig{fig:t1-1s}{周期增大后，同一实际时长内完成的振动次数减少；周期差夸张绘制}{\begin{tikzpicture}[x=.9cm,y=.8cm]
\draw[->](0,2.3)--(7.5,2.3)node[right]{$t$};\draw[->](0,0)--(7.5,0)node[right]{$t$};
\draw[blue,thick,domain=0:6.8,samples=200]plot(\x,{2.3+.6*cos(180*\x)});
\draw[red,thick,domain=0:6.8,samples=200]plot(\x,{.6*cos(138.4615*\x)});
\draw[dashed](6.5,-.8)--(6.5,3.3);\node[above]at(6.5,3.3){同一实际时刻};
\draw[<->](0,3.2)--(2,3.2)node[midway,above]{$P_0$};\draw[<->](0,-1)--(2.6,-1)node[midway,below]{$P>P_0$};
\node[left]at(0,2.3){原温};\node[left]at(0,0){升温后};\end{tikzpicture}}



\par\medskip\Needspace{4\baselineskip}\noindent\textbf{由实际时间与钟面读数之差求日差并检查走慢}\par\nobreak
实际一天取 $t=86400\,\mathrm s$，故少走的时间为
\[\Delta t=t-t_{\text{钟}}=t\left(1-\frac{P_0}{P}\right)
=86400\left(1-\frac1{\sqrt{1.00030}}\right)
\simeq\boxed{12.96\,\mathrm s}.\]
图中摆长增加，所以每次摆动用时增加，同一实际时间内计数减少，结果应是走慢。检查一阶近似时，令$\varepsilon=\alpha\Delta T=3.0\times10^{-4}$，分别对周期与计数因子使用幂函数展开，
\[
\frac P{P_0}=(1+\varepsilon)^{1/2}
=1+\frac{\varepsilon}{2}+O(\varepsilon^2),\qquad
\frac{P_0}{P}=(1+\varepsilon)^{-1/2}
=1-\frac{\varepsilon}{2}+O(\varepsilon^2).
\]
因此$\Delta t/t=1-P_0/P\simeq\varepsilon/2$，代入得$\Delta t\simeq86400\times(3.0\times10^{-4})/2=12.96\,\mathrm s$。这与直接计算相符，也表明日差的相对量级约为$1.5\times10^{-4}$。
\end{solution}
\par\medskip\noindent\begin{minipage}{\textwidth}
\begin{example}[受边界约束的肥皂膜细线]\label{t1-2}
两条光滑水平平行边相距 $d$，不可伸长的轻细线两端固定在左侧 $A,B$，长度 $\ell>\pi d/2$。细线两侧最初有肥皂膜，刺破左侧膜，右侧膜保留，表面张力系数为 $\sigma$。允许细线贴着光滑边界，求平衡形状和张力。见图\ref{fig:t1-2q}。
\end{example}
\fig{fig:t1-2q}{刺破左侧后保留右侧双面膜；图示细线尚未达到平衡}{\begin{tikzpicture}

\fill[blue!8](0,2)--(5,2)--(5,0)--(0,0)..controls(2.5,0)and(2.5,2)..(0,2)--cycle;
\draw[world](-.6,0)--(5,0) (-.6,2)--(5,2);\draw[very thick](0,2)..controls(2.5,2)and(2.5,0)..(0,0);\fill(0,2)circle(.05)node[left]{$A$};\fill(0,0)circle(.05)node[left]{$B$};\node at(.7,1){已刺破};\node at(3.7,1){右侧双面膜};\draw[dim](5.5,0)--(5.5,2)node[midway,right]{$d$};

\end{tikzpicture}}
\end{minipage}
\begin{solution}\textbf{在自由弧上用微元平衡确定恒定曲率}\par\nobreak
刺破左侧膜后，细线自由段只受到右侧双面膜的分布拉力。先不猜整条线的形状，在自由段取一小段 $\dd s$：两端张力大小相等，切向分量相互抵消；法向合力为 $\tau\dd\varphi$。两面膜对这一段的合力为 $2\sigma\dd s$，因此
\[\tau\dd\varphi=2\sigma\dd s,\qquad \frac1r=\frac{\dd\varphi}{\dd s}=\frac{2\sigma}{\tau}.\]
张力沿光滑轻线不变，自由段的曲率也不变，因而是圆弧。



\par\medskip\Needspace{4\baselineskip}\noindent\textbf{用上下边界筛选相切半圆与贴边段}\par\nobreak
还须使用图\ref{fig:t1-2q}的上下边界：细线不能越过两条水平边。由于 $\ell>\pi d/2$，连接上下边界的半圆已经不能容纳全部线长，平衡形状应包含贴边直线段。
在图\ref{fig:t1-2s}中，圆弧与边界光滑相切于 $C,D$，两点间竖直距离为 $d$，所以 $CD$ 为右半圆弧，$r=d/2$。\fig{fig:t1-2s}{受边界约束的肥皂膜细线：解答过程}{\begin{tikzpicture}[x=1cm,y=1cm]
\fill[blue!8] (1.8,2)--(5,2)--(5,0)--(1.8,0) arc[start angle=-90,end angle=90,radius=1]--cycle;\draw[world] (0,0)--(5,0) (0,2)--(5,2);\draw[very thick] (0,2)--(1.8,2) arc[start angle=90,end angle=-90,radius=1]--(0,0);\node[left] at (0,2){$A$};\node[left] at (0,0){$B$};\node[above] at (1.8,2){$C$};\node[below] at (1.8,0){$D$};\draw[guide] (1.8,1)--(2.8,1);\node at (1.8,1){$O$};\node[below] at (2.3,1){$r$};\draw[force] (2.8,1)--(4,1) node[right]{$2\sigma\,\dd s$};\node at (1,-.8){$AC=BD$};\draw[force] (1.8,2)--(.8,2) node[midway,above]{$\tau$};\draw[force] (1.8,0)--(.8,0) node[midway,below]{$\tau$};
\end{tikzpicture}}



\par\medskip\Needspace{4\baselineskip}\noindent\textbf{由半径求张力，再由总线长求贴边长度}\par\nobreak
由微元平衡式立即得到
\[\boxed{\tau=2\sigma r=\sigma d}.\]
圆弧长为 $\pi d/2$。其余部分由 $AC,DB$ 分担；上下对称且圆弧两端位于同一竖线上，故
\[\boxed{AC=BD=\frac{\ell-\pi d/2}{2}}.\]


\par\medskip\Needspace{4\baselineskip}\noindent\textbf{按真实膜区域检查整体受力与线长变化}\par\nobreak
图中蓝色膜仅在圆弧右侧，贴着上下边界的两段直线内侧已经没有膜，不能再给它们画出持续的横向膜力。再取整个半圆弧验算：膜力的水平合力为 $2\sigma d$ 向右，$C,D$ 处各有一个向左的张力，故 $2\tau=2\sigma d$，与局部推导一致。线长增大只使两段贴边线等量变长，半径和张力均保持不变。
\end{solution}
\par\medskip\noindent\begin{minipage}{\textwidth}
\begin{example}[双U形管的液面联动]\label{t1-3}
四根等截面竖管 $a,b,c,d$ 构成双U形装置：$a,b$ 下端由水银连通，$c,d$ 下端由水银连通，$b,c$ 顶端封有同一团气体，$a,d$ 开口。温度降低后 $c$ 中水银上升 $\Delta h$。求其余液面位移与气体压强变化。见图\ref{fig:t1-3q}，大气压恒定。
\end{example}
\fig{fig:t1-3q}{双U形管的液面联动：题设}{\begin{tikzpicture}[x=1cm,y=1cm]
\draw[world] (0,3)--(0,0)--(1.1,0)--(1.1,2.6)--(2.2,2.6)--(2.2,0)--(3.3,0)--(3.3,3);\draw[world] (.3,3)--(.3,.3)--(.8,.3)--(.8,2.9)--(2.5,2.9)--(2.5,.3)--(3,.3)--(3,3);\foreach \x/\t in {.15/a,.95/b,2.35/c,3.15/d}{\node[above] at (\x,3.2){$\t$};}\draw[thick] (0,2)--(.3,2) (.8,1.4)--(1.1,1.4) (2.2,1.4)--(2.5,1.4) (3,2)--(3.3,2);\node at (1.65,2.3){气体};
\end{tikzpicture}}
\end{minipage}
\begin{solution}\textbf{对两组独立水银分别列体积守恒}\par\nobreak
以竖直向上为正，设四液面位移为 $\delta h_a,\delta h_b,\delta h_c,\delta h_d$。图\ref{fig:t1-3q}中左、右两组水银彼此独立，每组的水银量都不变；又因四管等截面，分别有
\[\delta h_a+\delta h_b=0,\qquad \delta h_c+\delta h_d=0.\]


\par\medskip\Needspace{4\baselineskip}\noindent\textbf{以同一气体压强联系两组液面高度差}\par\nobreak
两外侧液面承受同一大气压，两内侧液面承受同一团气体的压强。沿左、右水银各从开口面走到气体面，得到
\[p=p_0+\rho g(h_a-h_b)=p_0+\rho g(h_d-h_c).\]
因此两组高度差的变化相等，即 $\delta h_a-\delta h_b=\delta h_d-\delta h_c$。题给 $\delta h_c=\Delta h$，依次代入上述关系，便确定图\ref{fig:t1-3s}中的四个箭头。\fig{fig:t1-3s}{双U形管的液面联动：解答过程}{\begin{tikzpicture}[x=1cm,y=1cm]
\foreach \shift/\title in {0/初态,5.5/末态}{\begin{scope}[xshift=\shift cm]
\draw[world] (0,3)--(0,0)--(1.1,0)--(1.1,2.6)--(2.2,2.6)--(2.2,0)--(3.3,0)--(3.3,3);
\draw[world] (.3,3)--(.3,.3)--(.8,.3)--(.8,2.9)--(2.5,2.9)--(2.5,.3)--(3,.3)--(3,3);
\foreach \x/\t in {.15/a,.95/b,2.35/c,3.15/d}{\node[above] at (\x,3.2){$\t$};}
\node at (1.65,3.9){同一团气体};\draw[guide] (1.65,3.6)--(1.65,2.75);\node at (1.65,-.5){\title};
\draw[guide] (-.2,0)--(3.6,0);
\end{scope}}
\draw[thick] (0,2)--(.3,2) (.8,1.4)--(1.1,1.4) (2.2,1.4)--(2.5,1.4) (3,2)--(3.3,2);
\draw[dim] (3.7,1.4)--(3.7,2) node[midway,right]{$H_i$};
\begin{scope}[xshift=5.5cm]
\draw[thick,blue] (0,1.8)--(.3,1.8) (.8,1.6)--(1.1,1.6) (2.2,1.6)--(2.5,1.6) (3,1.8)--(3.3,1.8);
\draw[guide] (-.1,2)--(.5,2) (.6,1.4)--(1.3,1.4) (2,1.4)--(2.7,1.4) (2.8,2)--(3.5,2);
\foreach \x/\a/\b in {.15/2/1.8,.95/1.4/1.6,2.35/1.4/1.6,3.15/2/1.8}{\draw[force,blue] (\x,\a)--(\x,\b);}
\draw[thick] (3.7,1.6)--(3.7,1.8);\draw (3.6,1.6)--(3.8,1.6) (3.6,1.8)--(3.8,1.8);\node[right] at (3.85,1.7){$H_f$};
\end{scope}
\node at (4.45,-1.1){$H_i=h_a-h_b=h_d-h_c,\qquad H_f=H_i-2\Delta h$};
\end{tikzpicture}}



\par\medskip\Needspace{4\baselineskip}\noindent\textbf{解出四液面位移，再求压强增量}\par\nobreak
结果为
\[\boxed{\delta h_b=\delta h_c=\Delta h,\qquad \delta h_a=\delta h_d=-\Delta h}.\]
图\ref{fig:t1-3s}中虚线保留各液面的初始位置，蓝色短线表示末位置；两幅图共用同一零高度。左组气体侧液面升高 $\Delta h$，大气侧液面降低 $\Delta h$，所以 $H_f=(h_a-\Delta h)-(h_b+\Delta h)=H_i-2\Delta h$。右组也给出同样变化，故
\[\boxed{p_f-p_i=-2\rho g\Delta h}.\]
这个倍数来自两边液面同时移动，不能只取内侧的一次位移。结论也符合降温后的趋势：气体压强降低，外侧大气把水银推入气体侧。整个推导只用了液柱平衡和水银体积守恒；若还要求末温，才需要补充气体体积和状态方程。
\end{solution}
\par\medskip\noindent\begin{minipage}{\textwidth}
\begin{example}[漏气水银气压计]\label{t1-4}
气压计顶部有少量空气。大气压为 $76\,\mathrm{cmHg}$ 时读数为 $74\,\mathrm{cm}$，顶部空气柱长 $8\,\mathrm{cm}$；后来读数为 $72\,\mathrm{cm}$，温度不变。槽面高度不变，求后来的真实大气压。见图\ref{fig:t1-4q}。
\end{example}
\fig{fig:t1-4q}{漏气水银气压计：题设}{\begin{tikzpicture}[x=1cm,y=1cm]
\draw[world] (0,0)--(0,4.1)--(1,4.1)--(1,0);\fill[gray!25] (.02,0) rectangle (.98,3.7);\draw[guide] (-.5,0)--(2,0);\draw[dim] (1.5,0)--(1.5,3.7) node[midway,right]{$74\,\mathrm{cm}$};\draw[dim] (-.5,3.7)--(-.5,4.1) node[midway,left]{$8$};\node at (.5,4.6){管顶固定};
\end{tikzpicture}}
\end{minipage}
\begin{solution}\textbf{用初读数求顶部空气压强与固定管顶高度}\par\nobreak
先从初态确定顶部空气的压强与管顶位置。图\ref{fig:t1-4q}中槽面作为固定零高度，管顶高度为 $H=74+8=82\,\mathrm{cm}$。水银柱只承担两端压差，所以初态满足
\[p_{01}=p_{a1}+74\,\mathrm{cmHg},\qquad p_{a1}=2\,\mathrm{cmHg}.\]


\par\medskip\Needspace{4\baselineskip}\noindent\textbf{由固定几何求第二次空气柱长度}\par\nobreak
读数变为 $72\,\mathrm{cm}$ 时，管顶没有移动，故空气柱不是原来的8cm，而是 $\ell_2=82-72=10\,\mathrm{cm}$。图\ref{fig:t1-4s}把两次空气柱长度及其对应压强并列表示。\fig{fig:t1-4s}{管顶和槽面固定，读数降低使顶部空气柱增长}{\begin{tikzpicture}[x=1cm,y=1.15cm]
\foreach \a/\hh/\lab/\air in {0/3.7/{74\,\mathrm{cm}}/{8\,\mathrm{cm}},4.5/3.6/{72\,\mathrm{cm}}/{10\,\mathrm{cm}}}{\begin{scope}[xshift=\a cm]
\draw[thick](0,-.2)--(0,4.1)--(.6,4.1)--(.6,-.2);\fill[gray!35](0,0)rectangle(.6,\hh);\fill[blue!8](0,\hh)rectangle(.6,4.1);
\draw[dashed](-.7,0)--(1,0);\draw[<->](1,0)--(1,\hh)node[midway,right]{$\lab$};
\draw[<->](-.6,\hh)--(-.6,4.1)node[midway,left]{$\air$};\node[above]at(.3,4.3){固定管顶};\end{scope}}
\node at(2.5,-.8){$p_{01}=p_{a1}+74,\qquad p_{02}=p_{a2}+72$（cmHg）};\end{tikzpicture}}



\par\medskip\Needspace{4\baselineskip}\noindent\textbf{用顶部空气的等温关系求真实大气压}\par\nobreak
顶部空气量不变、温度不变，且管截面相同，因而 $p_a\ell$ 不变。取压强单位为cmHg、长度单位为cm，有
\[p_{a2}\times10=2\times8,\qquad p_{a2}=1.6\,\mathrm{cmHg}.\]
由第二次水银柱的平衡关系，真实大气压等于顶部空气压强与水银柱压强之和：
\[\boxed{p_{02}=p_{a2}+72=73.6\,\mathrm{cmHg}}.\]
空气柱变长、空气压强减小，与等温膨胀一致。若把72cm直接作为大气压，会漏掉顶部空气承担的1.6cmHg；若令空气柱仍长8cm，则违背图中固定管顶、固定槽面的几何条件。
\end{solution}
\par\medskip\noindent\begin{minipage}{\textwidth}
\begin{example}[双挡板之间的加热]\label{t1-5}
水平气缸内轻活塞初始靠左挡板 $B$，气体体积 $V_0$、压强 $0.9p_0$、温度297K，外压 $p_0$。右挡板 $A$ 对应体积 $1.1V_0$。缓慢加热到399.3K，求离开 $B$、到达 $A$ 的温度和末压，画出 $p$--$V$ 路径。见图\ref{fig:t1-5q}。
\end{example}
\fig{fig:t1-5q}{双挡板之间的加热：题设}{\begin{tikzpicture}[x=1cm,y=1cm]
\draw[world] (0,2)--(0,0)--(4.7,0) (0,2)--(4.7,2);\fill[blue!6] (.02,.02) rectangle (3,1.98);\fill[pattern=north east lines] (3,0) rectangle (3.13,2);\node at (1.5,1){$0.9p_0,V_0,297\,\mathrm K$};\node at (4.1,1){$p_0$};\node[above] at (3,2.2){$B$};\draw[guide] (3.6,0)--(3.6,2);\node[above] at (3.6,2.2){$A$};
\end{tikzpicture}}
\end{minipage}
\begin{solution}\textbf{活塞靠左挡板时，以支持力归零确定离开温度}\par\nobreak
初态气压小于外压，活塞被大气压在左挡板 $B$ 上。加热初期，$B$ 提供向右支持力 $N_B=(p_0-p)S$，只要 $p<p_0$，活塞就不能自由向右运动，体积保持 $V_0$。当 $N_B$ 首次减为零，$p_2=p_0$，气体由等容关系给出
\[\boxed{T_2=297\frac{p_0}{0.9p_0}=330\,\mathrm K}.\]


\par\medskip\Needspace{4\baselineskip}\noindent\textbf{活塞自由运动时，以等压膨胀确定接触右挡板温度}\par\nobreak
离开 $B$ 后，轻活塞无摩擦地缓慢移动，不再有挡板支持力，故两侧气压相等。此段等压，直至体积由 $V_0$ 增大到 $1.1V_0$，有
\[\boxed{T_3=T_2\frac{1.1V_0}{V_0}=363\,\mathrm K}.\]
图\ref{fig:t1-5s}中的点2、3分别对应“刚离开 $B$”和“刚接触 $A$”；水平段 $2\to3$ 才是自由运动段。\fig{fig:t1-5s}{双挡板之间的加热：解答过程}{\begin{tikzpicture}[x=1cm,y=1cm]
\draw[axis] (0,0)--(4.8,0) node[right]{$V/V_0$};\draw[axis] (0,0)--(0,3.4) node[above]{$p/p_0$};\draw (.15,-.08)--(.25,.08)--(.35,-.08)--(.45,.08);\draw (-.08,.15)--(.08,.25)--(-.08,.35)--(.08,.45);\draw[ray] (1,.6)--(1,1.6);\draw[ray] (1,1.6)--(3,1.6);\draw[ray] (3,1.6)--(3,2.6);\foreach \x/\y/\t in {1/.6/1,1/1.6/2,3/1.6/3,3/2.6/4}{\fill (\x,\y) circle(.04);\node[above left] at (\x,\y){$\t$};}\node[below] at (1,0){$1$};\node[below] at (3,0){$1.1$};\node[left] at (0,.6){$0.9$};\node[left] at (0,1.6){$1$};\node[left] at (0,2.6){$1.1$};\node at (2,-.8){坐标轴截取局部};
\end{tikzpicture}}



\par\medskip\Needspace{4\baselineskip}\noindent\textbf{接触右挡板后，以等容升压求末压}\par\nobreak
由于 $399.3\,\mathrm K>363\,\mathrm K$，末态已在右挡板上。继续加热时体积固定为 $1.1V_0$，于是
\[\boxed{p_4=p_3\frac{T_4}{T_3}=p_0\frac{399.3}{363}=1.1p_0}.\]


\par\medskip\Needspace{4\baselineskip}\noindent\textbf{核对两挡板的反力方向与非负条件}\par\nobreak
图中 $1\to2$、$3\to4$ 两条竖线分别是两个挡板建立的等容约束。后一段右挡板向左的支持力为 $(p-p_0)S$，从零增加到 $0.1p_0S$，始终非负；前一段左挡板的支持力则由 $0.1p_0S$ 减到零。两段力的方向均与边界阻挡方向一致，因此这一分段确实可实现。
\end{solution}
\par\medskip\noindent\begin{minipage}{\textwidth}
\begin{example}[沙袋活塞与温度反推]\label{t1-6}
轻活塞压着气体并承受沙袋，初始靠在支环上，气柱高 $H_0$、压强 $p_0$。加热60K时活塞恰能离环，再继续加热使其上升到 $1.5H_0$。保持此时温度，逐渐移去全部沙袋，气柱最终高 $1.8H_0$。外压恒为 $p_0$，求初温及末温。见图\ref{fig:t1-6q}。
\end{example}
\fig{fig:t1-6q}{沙袋活塞与温度反推：题设}{\begin{tikzpicture}[x=1cm,y=1cm]
\draw[world] (0,3)--(0,0)--(2,0)--(2,3);\fill[pattern=north east lines] (0,1.2) rectangle (2,1.35);\draw[world] (.6,1.35) rectangle (1.4,1.8);\node at (1,1.6){沙};\node at (1,.55){$p_0,T_0$};\draw[dim] (2.4,0)--(2.4,1.2) node[midway,right]{$H_0$};\node[left] at (0,1.2){支环};
\end{tikzpicture}}
\end{minipage}
\begin{solution}\textbf{区分支环承托状态与带沙袋的自由平衡压强}\par\nobreak
设沙袋质量为 $m$、活塞面积为 $S$。只要活塞带着全部沙袋自由上升，平衡压强就是常量 $P=p_0+mg/S$。图\ref{fig:t1-6q}的初态尚由支环承托，故初始气压 $p_0$ 不必等于 $P$。


\par\medskip\Needspace{4\baselineskip}\noindent\textbf{利用移沙等温段反求带载压强}\par\nobreak
从末两态反推 $P$ 最直接：移沙时温度保持 $T_1$，活塞下气体量不变，而体积正比于气柱高度。按图\ref{fig:t1-6s}中的等温段 $C\to D$，
\[P\,(1.5SH_0)=p_0(1.8SH_0),\qquad P=1.2p_0.\]


\par\medskip\Needspace{4\baselineskip}\noindent\textbf{用离环前等容关系反求初温}\par\nobreak
再回到最初的加热阶段。活塞离环前高度仍是 $H_0$；加热60K后支环支持力恰降至零，此时才达到自由平衡压强 $P$。\fig{fig:t1-6s}{局部放大的状态路径，轴上数字保留实际无量纲取值}{\begin{tikzpicture}[x=1.1cm,y=1cm]
\draw[->](-.4,0)--(5.3,0)node[right]{$V/(SH_0)$};\draw[->](0,-.5)--(0,2.7)node[above]{$p/p_0$};
\draw[->,blue,thick](0,0)--(0,1.6);\draw[->,red,thick](0,1.6)--(2.5,1.6);
\draw[->,orange,thick,domain=2.5:4,samples=60]plot(\x,{8*(1.8/(1+\x/5)-1)});
\fill(0,0)circle(1.5pt)node[above right]{$A$};\fill(0,1.6)circle(1.5pt)node[above right]{$B$};
\fill(2.5,1.6)circle(1.5pt)node[above]{$C$};\fill(4,0)circle(1.5pt)node[above right]{$D$};
\node[left]at(0,0){$1$};\node[left]at(0,1.6){$1.2$};\node[below]at(0,-.4){$1$};\node[below]at(2.5,-.1){$1.5$};\node[below]at(4,-.1){$1.8$};
\node at(2.5,-1.2){$A\to B$等容；$B\to C$等压；$C\to D$等温};\end{tikzpicture}}

等容状态关系给出
\[\frac{T_0+60}{T_0}=\frac P{p_0}=1.2,\qquad \boxed{T_0=300\,\mathrm K}.\]


\par\medskip\Needspace{4\baselineskip}\noindent\textbf{用带载等压关系求最终温度并检查移沙阶段}\par\nobreak
因此离环温度为360K。此后保留全部沙袋继续加热，气体等压膨胀到原高度的1.5倍：
\[\boxed{T_1=360\frac{1.5H_0}{H_0}=540\,\mathrm K}.\]
逐渐移沙是另一个等温膨胀过程，末温仍为540K；高度比1.8/1.5通过压强下降来满足状态方程，不能再次用来提高温度。三阶段分别保持体积、压强和温度，图中同一高度或同一温度的对应关系正是反推未知初温的依据。
\end{solution}
\par\medskip\noindent\begin{minipage}{\textwidth}
\begin{example}[提起气缸的临界压强]\label{t1-7}
质量 $M$ 的绝热竖直气缸放在地面，截面 $A$。轻活塞下气体初态为 $p_0,T_0$，外压为 $p_0$。气缸由小支脚承托，缸底外表面始终与大气相通；忽略壁厚、活塞摩擦及气体质量，气缸足够长且活塞始终密封。缓慢上提活塞，直到气缸恰离地。气体满足 $pV^\gamma=\text{常量}$。求离地温度及必要条件。见图\ref{fig:t1-7q}。
\end{example}
\fig{fig:t1-7q}{上提活塞使气体膨胀；缸底由支脚架空，大气能够作用于底板外侧}{\begin{tikzpicture}

\draw[world](0,3)--(0,0)--(2,0)--(2,3);\fill[blue!8](0,0)rectangle(2,1.3);\draw[very thick](0,1.3)--(2,1.3);\draw[force](1,1.3)--(1,2.7)node[above]{$F$};\node at(1,.6){$p,T$};\node[right]at(2,2){$p_0$};
\draw[world](.15,0)--(.15,-.4) (1.85,0)--(1.85,-.4) (-.4,-.4)--(2.4,-.4);\draw[flow](1,-.9)--(1,-.1)node[left]{$p_0$};\node[right]at(2.4,-.2){底部通大气};

\end{tikzpicture}}
\end{minipage}
\begin{solution}\textbf{以气缸壳体为对象，用地面反力归零求临界压强}\par\nobreak
应把气缸壳体与气体分别作为研究对象。上提轻活塞会使气体膨胀、压强下降；真正使壳体离地的是底板内外的压差。图\ref{fig:t1-7s}中底板外侧大气压向上作用，内侧气压向下作用。壳体竖直平衡为
\[N+p_0A-pA-Mg=0.\]
刚离地时地面对壳体的支持力 $N=0$，故临界压强为
\[\boxed{p_f=p_0-\frac{Mg}{A}}.\]
这一步只依赖壳体受力，与气体的绝热指数无关。\fig{fig:t1-7s}{提起气缸的临界压强：解答过程}{\begin{tikzpicture}[x=1cm,y=1cm]
\draw[world] (0,2)--(0,0)--(2,0)--(2,2);\draw[force] (.5,-.9)--(.5,-.05) node[above left]{$p_0A$};\draw[force] (1.3,.05)--(1.3,-.9) node[below]{$p_f A$};\draw[force] (2,1)--(2,-.1) node[right]{$Mg$};\node at (1,2.6){壳体临界受力：$N=0$};
\end{tikzpicture}}



\par\medskip\Needspace{4\baselineskip}\noindent\textbf{由准静态绝热路径把临界压强转换为温度}\par\nobreak
再对封闭气体使用题给的准静态绝热关系。由 $pV^\gamma=p_0V_0^\gamma$ 得 $V_f/V_0=(p_0/p_f)^{1/\gamma}$；结合状态方程，
\[\frac{T_f}{T_0}=\frac{p_fV_f}{p_0V_0}
=\left(\frac{p_f}{p_0}\right)^{(\gamma-1)/\gamma}.\]
代入临界压强，得
\[\boxed{T_f=T_0\left(1-\frac{Mg}{p_0A}\right)^{(\gamma-1)/\gamma}}.\]


\par\medskip\Needspace{4\baselineskip}\noindent\textbf{检查正压条件与无限膨胀极限}\par\nobreak
有限体积的气体压强必须为正，因此需要 $\boxed{p_0A>Mg}$。等号只对应无限膨胀、压强趋零的极限，不能在有限位移处离地；若 $p_0A<Mg$，即便内部接近真空，大气在底板上所能提供的最大向上作用也不足以克服重力。所得 $T_f<T_0$ 则与膨胀对外做功、内能降低相符。
\end{solution}
\par\medskip\noindent\begin{minipage}{\textwidth}
\begin{example}[三截面管与不可伸长细线]\label{t1-8}
水平管的三个区段面积依次为 $2S,S/2,S$，中段长 $2\ell$。两轻活塞 $A,B$ 分别在左右宽段，距相邻台阶均为 $\ell$，用长 $4\ell$ 的轻绳连接。初态 $p_0,T_0$，外压 $p_0$。缓慢升温到 $T\ge T_0$，求压强的分段表达式。活塞不能通过较窄的台阶。见图\ref{fig:t1-8q}。
\end{example}
\fig{fig:t1-8q}{两活塞间气体占据三个截面不同的区段；各轴向长度按同一比例画出}{\begin{tikzpicture}

\draw[world](-.7,1.6)--(2,1.6)--(2,.4)--(4,.4)--(4,.8)--(6.5,.8) (-.7,-1.6)--(2,-1.6)--(2,-.4)--(4,-.4)--(4,-.8)--(6.5,-.8);
\draw[very thick](1,-1.6)--(1,1.6) (5,-.8)--(5,.8);\draw[world](1,0)--(5,0);\node[above]at(1,1.6){$A$};\node[above]at(5,.8){$B$};\node at(.1,.6){$2S$};\node at(3,.85){$S/2$};\node at(6,.3){$S$};
\draw[dim](1,-2.1)--(2,-2.1)node[midway,below]{$\ell$};\draw[dim](2,-2.1)--(4,-2.1)node[midway,below]{$2\ell$};\draw[dim](4,-2.1)--(5,-2.1)node[midway,below]{$\ell$};\node at(3,2.5){轻绳长$4\ell$，中段长度$2\ell$};

\end{tikzpicture}}
\end{minipage}
\begin{solution}\textbf{按截面分段求初始体积}\par\nobreak
图\ref{fig:t1-8q}中两活塞之间的初体积要逐段计算：
\[V_0=2S\ell+\frac S2(2\ell)+S\ell=4S\ell.\]


\par\medskip\Needspace{4\baselineskip}\noindent\textbf{分别列自由活塞平衡并消去同一绳张力}\par\nobreak
先考察两活塞都未接触台阶的阶段。若绳保持伸直，$A$ 上的压差力要求 $\tau=2S(p-p_0)$，$B$ 上则要求 $\tau=S(p-p_0)$。同一轻绳的张力相同，故只能有 $p=p_0,\tau=0$；这是准静态运动的极限，升温引起的体积增加决定移动方向。


\par\medskip\Needspace{4\baselineskip}\noindent\textbf{用共同位移求体积，确定首先接触的活塞}\par\nobreak
取两活塞共同向左移动的距离为 $x$。

左侧宽段长度增加 $x$，右侧宽段长度减少 $x$，中段不变，因此
\[V(x)=2S(\ell+x)+S\ell+S(\ell-x)=4S\ell+Sx.\]
可见向左移动反而使体积增大。升温时正沿此方向运动；首先到达台阶的是 $B$，其临界位移 $x=\ell$，对应 $V_*=5S\ell$。\fig{fig:t1-8s}{共同左移使左侧增容超过右侧减容；$B$到达右台阶后两活塞均固定}{\begin{tikzpicture}

\foreach \xx/\lab in {0/初态,1/临界态}{\begin{scope}[yshift={-3.6*\xx*1cm}]
\draw[world](-.7,1.3)--(2,1.3)--(2,.4)--(4,.4)--(4,.8)--(6,.8) (-.7,-1.3)--(2,-1.3)--(2,-.4)--(4,-.4)--(4,-.8)--(6,-.8);\draw[very thick]({1-\xx},-1.3)--({1-\xx},1.3) ({5-\xx},-.8)--({5-\xx},.8);\draw[world]({1-\xx},0)--({5-\xx},0);\node[above]at({1-\xx},1.3){$A$};\node[above]at({5-\xx},.8){$B$};\node[right]at(6.2,0){\lab};\end{scope}}
\draw[flow](1,2)--(0,2)node[midway,above]{$x$};\draw[flow](5,2)--(4,2)node[midway,above]{$x$};\node at(3,-5.6){$V=4S\ell+Sx,\quad 0\le x\le\ell$};

\end{tikzpicture}}



\par\medskip\Needspace{4\baselineskip}\noindent\textbf{在接触事件处衔接等压和等容两段}\par\nobreak
图\ref{fig:t1-8s}中共同左移至$x=\ell$时压强仍为$p_0$，所以
\[T_*=T_0\frac{V_*}{V_0}=\frac54T_0.\]
到达图\ref{fig:t1-8s}下部的位置后，台阶阻止 $B$ 继续左移；若 $A$ 再向左，绳长便会增加，因此 $A$ 也固定。继续加热属于等容过程，最终得到
\[\boxed{\frac p{p_0}=\begin{cases}1,&1\le T/T_0\le5/4,\\[1mm]4T/(5T_0),&T/T_0\ge5/4.\end{cases}}\]
\fig{fig:h1:e8path}{共同移动段保持$p=p_0$；$B$接触台阶后沿等容直线升压}{\begin{tikzpicture}

\draw[axis](0,0)--(5.8,0)node[right]{$T/T_0$};\draw[axis](0,0)--(0,3)node[above]{$p/p_0$};\draw[flow,blue](3,1.5)--(3.75,1.5);\draw[flow,red](3.75,1.5)--(5.4,2.16);\draw[guide](3,0)--(3,1.5) (3.75,0)--(3.75,1.5)--(0,1.5);\node[below]at(3,0){$1$};\node[below]at(3.75,0){$5/4$};\node[left]at(0,1.5){$1$};\node at(1.4,2.4){共同左移};\node at(4.5,2.8){接触后等容};

\end{tikzpicture}}
图\ref{fig:h1:e8path}的两段在$(T/T_0,p/p_0)=(5/4,1)$连续衔接。接触后体积不变，所以$p/T$为常量，这条红线的反向延长线通过原点；它与蓝色等压段代表不同约束。


\fig{fig:h1:e8forces}{图取接触后的$p>p_0$；压力已合并成内外压差力，台阶只在$B$上提供向右反力}{\begin{tikzpicture}

\begin{scope}\draw[world](0,0)rectangle(1,.6);\node at(.5,.3){$A$};\draw[force](0,.3)--(-1.5,.3)node[left]{$2S(p-p_0)$};\draw[force](1,.3)--(2.4,.3)node[right]{$\tau$};\end{scope}
\begin{scope}[yshift=-2.8cm]\draw[world](0,0)rectangle(1,.6);\node at(.5,.3){$B$};\draw[force](0,.3)--(-1.5,.3)node[left]{$\tau$};\draw[force](1,.45)--(2.4,.45)node[right]{$S(p-p_0)$};\draw[force](1,.1)--(2.4,.1)node[right]{$N_B$};\end{scope}

\end{tikzpicture}}
图\ref{fig:h1:e8forces}将两个活塞分开，绳对$A$向右、对$B$向左。台阶接触前没有$N_B$，两个平衡式只能共同给出$p=p_0,\tau=0$；接触后多出一个允许的向右反力，才可维持$p>p_0$的静止状态。




\par\medskip\Needspace{4\baselineskip}\noindent\textbf{核对绳张力与台阶反力均非负}\par\nobreak
还需检验后段受力。$A$ 平衡要求 $\tau=2S(p-p_0)\ge0$；在 $B$ 上，绳向左拉，压差力和台阶支持力向右，因此 $N_B=\tau-S(p-p_0)=S(p-p_0)\ge0$。图中 $N_B$ 的方向正能阻止 $B$ 穿过台阶。两段在 $T_*$ 处压强连续，张力与支持力均由零开始增加。
\end{solution}
\par\medskip\noindent\begin{minipage}{\textwidth}
\begin{example}[双活塞冷却的三段路径]\label{t1-9}
竖直台阶管下段面积 $2S=20\,\mathrm{cm^2}$，上段面积 $S=10\,\mathrm{cm^2}$。$A$ 质量1kg，初高 $L$；台阶高 $2L$；$B$ 质量0.5kg，初高 $4L$。两活塞用长 $3L$ 的绳相连，$B$ 还经定滑轮连接质量 $M$ 的砝码。初气压 $1.2\times10^5$Pa、温度800K，外压 $10^5$Pa，$g=10\,\mathrm{m/s^2}$。缓慢冷却到210K，求 $M$ 和各阶段状态。见图\ref{fig:t1-9q}。
\end{example}
\fig{fig:t1-9q}{两活塞间为独立轻绳；上方另一根绳从$B$绕过定滑轮连接外侧砝码}{\begin{tikzpicture}

\draw[world](0,0)--(0,2)--(.5,2)--(.5,4.7) (2,0)--(2,2)--(1.5,2)--(1.5,4.7);\fill[blue!8](0,1)rectangle(2,2);\fill[blue!8](.5,2)rectangle(1.5,4);
\draw[very thick](0,1)--(2,1) (.5,4)--(1.5,4);\draw[world](1,1)--(1,4);\draw[world](1,4)--(1,5.2)arc(180:0:.6)--(2.2,4.1);\draw[world](1.6,5.2)circle(.6);\fill(1.6,5.2)circle(.04);\draw[world](1.6,5.2)--(1.6,6.2);\draw[world](1.1,6.2)--(2.1,6.2);\draw[world](1.9,3.5)rectangle(2.5,4.1);\node at(2.2,3.8){$M$};
\node[left]at(0,1){$A:L$};\node[left]at(0,2){$2L$};\node[left]at(.5,4){$B:4L$};\node[left]at(.95,2.7){$\tau$};\node[right]at(2.2,4.6){$Mg$};\node at(1,-.6){气体在$A$上方、$B$下方};

\end{tikzpicture}}
\end{minipage}
\begin{solution}\textbf{分开画两活塞受力，利用初态反求砝码质量}\par\nobreak
先取竖直向上为正，并记 $\Delta p=p-p_0$。图\ref{fig:t1-9q}中气体位于 $A$ 上方、$B$ 下方，两个气体压力的方向相反。两活塞间绳张力为 $\tau$，上方经定滑轮的绳张力为 $Mg$，分别列平衡得
\[A:\quad\tau=10+2S\Delta p,\qquad
B:\quad\tau=S\Delta p+Mg-5.\]
\fig{fig:h1:e9forces}{初段两活塞均自由：气体净压力对$A$向下、对$B$向上，绳的两端拉力相反}{\begin{tikzpicture}

\begin{scope}\draw[world](0,0)rectangle(1,.6);\node at(.5,.3){$A$};\draw[force](.5,.6)--(.5,2)node[above]{$\tau$};\draw[force](.2,0)--(.2,-1.3)node[left]{$2S\Delta p$};\draw[force](.8,0)--(.8,-.8)node[right]{$m_Ag$};\end{scope}
\begin{scope}[xshift=5.5cm]\draw[world](0,0)rectangle(1,.6);\node at(.5,.3){$B$};\draw[force](.2,.6)--(.2,2)node[left]{$Mg$};\draw[force](.8,.6)--(.8,1.4)node[right]{$S\Delta p$};\draw[force](.2,0)--(.2,-1.3)node[left]{$\tau$};\draw[force](.8,0)--(.8,-.8)node[right]{$m_Bg$};\end{scope}

\end{tikzpicture}}
图\ref{fig:h1:e9forces}中$\Delta p=p-p_0$，只把内外压力合并，两个重力和两根绳的张力仍分别保留。若后续$A$接触台阶，须在其图中另加向下反力；两活塞间的绳松弛时，则应把两图中的$\tau$同时置零。


初态 $S=10^{-3}\,\mathrm{m^2}$、$\Delta p=2\times10^4\,\mathrm{Pa}$，所以 $\tau=50\,\mathrm N$，再由 $B$ 的方程得 $Mg=35\,\mathrm N$，即 $\boxed{M=3.5\,\mathrm{kg}}$。



\par\medskip\Needspace{4\baselineskip}\noindent\textbf{共同上移的等压冷却，直到下活塞接触台阶}\par\nobreak
两活塞同时自由且绳绷直时，联立平衡方程使 $p=1.2\times10^5\,\mathrm{Pa}$ 固定。设共同上移 $x$，则下段气体缩短 $x$、上段气体增长 $x$，
\[V=2S(L-x)+S(2L+x)=4SL-Sx.\]
因此冷却时两活塞上移。$A$ 首先在 $x=L$ 时接触台阶，此时 $z_A=2L,z_B=5L,V=3SL$，由等压关系得 $T=800\times3/4=600\,\mathrm K$。图\ref{fig:t1-9s}将这一步的装置位置与状态路径对应起来。\fig{fig:t1-9s}{双活塞冷却的三段路径：解答过程}{\begin{tikzpicture}[x=1cm,y=1cm]
\draw[axis] (0,0)--(5.4,0) node[right]{$V/(SL)$};\draw[axis] (0,0)--(0,3.5) node[above]{$p/(10^5\,\mathrm{Pa})$};\node[below left] at (0,0){$O$};\draw[ray] (4,2.4)--(3,2.4);\draw[ray] (3,2.4)--(3,1.4);\draw[ray] (3,1.4)--(1.8,1.4);\node[above] at (4,2.4){$800$K};\node[above left] at (3,2.4){$600$K};\node[right] at (3,1.4){$350$K};\node[below] at (1.8,1.4){$210$K};\node[left] at (0,2.4){$1.2$};\node[left] at (0,1.4){$0.7$};\foreach \x in {1.8,3,4}{\node[below] at (\x,0){$\x$};}\begin{scope}[shift={(-.5,-4)}]\foreach \x/\ha/\hb/\txt in {0/1/4/初态,2.5/2/5/接触台阶}{\begin{scope}[shift={(\x,0)},scale=.55]\draw[world] (0,0)--(0,2)--(.5,2)--(.5,5.5) (2,0)--(2,2)--(1.5,2)--(1.5,5.5);\draw[very thick] (0,\ha)--(2,\ha) (.5,\hb)--(1.5,\hb);\draw (1,\ha)--(1,\hb);\node[below] at (1,-.2){\txt};\end{scope}}\begin{scope}[shift={(5,0)},scale=.55]\draw[world] (0,0)--(0,2)--(.5,2)--(.5,5.5) (2,0)--(2,2)--(1.5,2)--(1.5,5.5);\draw[very thick] (0,2)--(2,2) (.5,3.8)--(1.5,3.8);\draw (1,2) .. controls (.3,2.4) and (1.7,3) .. (1,3.8);\node[below] at (1,-.2){绳松弛};\end{scope}\end{scope}
\end{tikzpicture}}



\par\medskip\Needspace{4\baselineskip}\noindent\textbf{固定体积冷却，直到两活塞间的绳张力归零}\par\nobreak
接触后，$A$不能继续上移；绳长限制又使$B$不能继续上移。此时$B$的自由平衡仍要求$\tau=S\Delta p+30$。只要此值为正，$B$就不能向下离开并使绳松弛后仍保持平衡，因而两活塞位置及体积暂时固定。冷却使$p$和$\tau$同时降低，到
\[\tau=0\quad\Longrightarrow\quad p=0.7\times10^5\,\mathrm{Pa}\]
时绳恰松弛，对应温度
\[T=800\frac{(0.7)(3)}{(1.2)(4)}=350\,\mathrm K.\]
在这一等容段，$A$ 上台阶向下的支持力为 $N=\tau-10-2S\Delta p=20-S\Delta p$，由0增至50N，方向始终可行。



\par\medskip\Needspace{4\baselineskip}\noindent\textbf{绳松弛后重新列上活塞平衡并求末位置}\par\nobreak
继续冷却时两活塞间绳已松弛。$B$ 由上绳拉力、重力和压差力平衡，使 $p$ 保持 $0.7\times10^5\,\mathrm{Pa}$；$A$ 则仍贴台阶。于是末体积与位置为
\[\boxed{V_f=3SL\frac{210}{350}=1.8SL,\quad z_{Af}=2L,\quad z_{Bf}=2L+\frac{V_f}{S}=3.8L}.\]
末态两活塞间距 $1.8L<3L$，绳确已松弛；$A$ 的支持力仍为50N。


\par\medskip\Needspace{4\baselineskip}\noindent\textbf{汇总关键状态并核对各段反力与几何边界}\par\nobreak
关键状态如下，其中 $\tau$ 指两活塞间绳张力。
\begin{center}\small\begin{tabular}{cccccc}\toprule $T/\mathrm K$&$p/(10^5\mathrm{Pa})$&$V/(SL)$&$z_A/L$&$z_B/L$&$\tau/\mathrm N$\\\midrule800&1.2&4&1&4&50\\600&1.2&3&2&5&50\\350&0.7&3&2&5&0\\210&0.7&1.8&2&3.8&0\\\bottomrule\end{tabular}\end{center}

\end{solution}

\section{本讲小结}
热膨胀问题应从基准长度出发，把材料的伸长传到所求的几何量或仪器读数。液面问题应区分一层界面与双面薄膜，并将局部曲率关系和整体边界条件共同使用。
对于液柱和活塞装置，气态方程、机械平衡与几何约束各承担不同作用：气态方程联系状态量，机械平衡确定压强，几何约束给出体积。先把这三类关系写清，再在接触、离开或张力降为零的位置切换条件，才能得到完整而可检验的过程。

\section{练习}
\par\medskip\noindent\begin{minipage}{\textwidth}
\begin{drill}[金属环孔的变化]\label{t1-10}
各向同性金属圆环在 $T_0$ 时内半径为 $a$、外半径为 $b$，在可自由膨胀的条件下均匀升温 $\Delta T$。线膨胀系数为 $\alpha$，求孔面积与环材料面积的一阶变化。见图\ref{fig:t1-10q}。
\end{drill}
\fig{fig:t1-10q}{金属环孔的变化：题设}{\begin{tikzpicture}[x=1cm,y=1cm]
\fill[gray!20,even odd rule] (0,0) circle(1.5) (0,0) circle(.7);\draw[world] (0,0) circle(1.5) (0,0) circle(.7);\draw[dim] (0,0)--(.7,0) node[midway,above]{$a$};\draw[dim] (0,0)--(0,-1.5) node[midway,left]{$b$};
\end{tikzpicture}}
\hyperref[sol:t1-10]{详解见第\pageref{sol:t1-10}页}。
\end{minipage}
\par\medskip\noindent\begin{minipage}{\textwidth}
\begin{drill}[润湿与毛细高度]\label{t1-11}
两根内半径分别为 $r,2r$ 的细管插入同一液体。液体密度 $\rho$、表面张力 $\sigma$，接触角均为 $60^\circ$。忽略弯月面体积，求液面高度及高度比。见图\ref{fig:t1-11q}。
\end{drill}
\fig{fig:t1-11q}{题设：两根管的半径比为$1:2$，两弯月面与管壁的液体内接触角同为$60^\circ$}{\begin{tikzpicture}

\draw[guide](-1,0)--(5,0)node[right]{槽面};\draw[world](0,-.5)--(0,2.6) (.5,-.5)--(.5,2.6) (3,-.5)--(3,2.6) (4,-.5)--(4,2.6);\draw[blue,thick](0,2.066987)arc(240:300:.5);\draw[blue,thick](3,1.133975)arc(240:300:1);
\node[above]at(.25,2.6){$r$};\node[above]at(3.5,2.6){$2r$};\draw[dim](1,0)--(1,2)node[midway,right]{$h_r$};\draw[dim](4.5,0)--(4.5,1)node[midway,right]{$h_{2r}$};\draw[guide](0,2.066987)--(.5,1.778312);\draw(0,1.866987)arc(-90:-30:.2);\node[left]at(-.1,1.85){$60^\circ$};

\end{tikzpicture}}
\hyperref[sol:t1-11]{详解见第\pageref{sol:t1-11}页}。
\end{minipage}
\par\medskip\noindent\begin{minipage}{\textwidth}
\begin{drill}[两室等温连通]\label{t1-12}
刚性容器两室体积分别为 $V,2V$，初压分别为 $3p_0,p_0$，内含同种理想气体。始末温度均为 $T_0$。打开阀门后求平衡压强，以及由第一室净流出的物质的量。见图\ref{fig:t1-12q}。
\end{drill}
\fig{fig:t1-12q}{两室等温连通：题设}{\begin{tikzpicture}[x=1cm,y=1cm]
\draw[world] (0,0) rectangle (2,2) (2.7,0) rectangle (6.7,2);\draw (2,1)--(2.7,1);\draw[world](2.2,.85)--(2.5,1.15)--(2.5,.85)--(2.2,1.15)--cycle;\node at (1,1){$V,3p_0$};\node at (4.7,1){$2V,p_0$};\node[above] at (2.35,1.2){$K$};
\end{tikzpicture}}
\hyperref[sol:t1-12]{详解见第\pageref{sol:t1-12}页}。
\end{minipage}
\par\medskip\noindent\begin{minipage}{\textwidth}
\begin{drill}[重活塞接触上挡板]\label{t1-13}
竖直活塞面积 $S$、质量 $m$，初始自由平衡，气柱高 $h_0$、温度 $T_0$，外压 $p_0$。上挡板在 $2h_0$。缓慢加热到 $3T_0$，求末压与挡板支持力。见图\ref{fig:t1-13q}。
\end{drill}
\fig{fig:t1-13q}{重活塞接触上挡板：题设}{\begin{tikzpicture}[x=1cm,y=1cm]
\draw[world] (0,3)--(0,0)--(2,0)--(2,3);\fill[pattern=north east lines] (0,1.2) rectangle (2,1.35);\draw[guide] (0,2.4)--(2,2.4);\draw[world,very thick](0,2.4)--(.3,2.4) (1.7,2.4)--(2,2.4);\node at(1,3.2){$m,\ S,\ p_0$};\node[right] at (2,1.2){$h_0$};\node[right] at (2,2.4){$2h_0$};\node at (1,.6){$T_0$};
\end{tikzpicture}}
\hyperref[sol:t1-13]{详解见第\pageref{sol:t1-13}页}。
\end{minipage}

\LectureAnswerSection

\begin{drillsolution}{t1-10}\textbf{用自由膨胀的位移场同时确定内外半径}\par\nobreak
自由、均匀加热时，圆环各点相对于中心 $O$ 的距离按同一比例变化。图\ref{fig:t1-10s}的虚线为初始边界、蓝线为加热后的边界，故
\[a'=a(1+\alpha\Delta T),\qquad b'=b(1+\alpha\Delta T).\]
孔的边界也是金属上的材料点，因而内半径同样增加。可以想象先用同种材料填孔，整体膨胀后再取出填充部分；这一操作不会改变原圆环边界的位置。\fig{fig:t1-10s}{金属环孔的变化：解答过程}{\begin{tikzpicture}[x=1cm,y=1cm]
\draw[guide] (0,0) circle(.7) (0,0) circle(1.5);\draw[world,blue] (0,0) circle(.84) (0,0) circle(1.8);\fill (0,0) circle(.03);\node[below left] at (0,0){$O$};\draw[flow] (1.5,0)--(1.8,0);\node[align=left] at (4,0){$a\to a(1+\alpha\Delta T)$\\$b\to b(1+\alpha\Delta T)$};
\end{tikzpicture}}



\par\medskip\Needspace{4\baselineskip}\noindent\textbf{先对内外半径平方，再统一舍去二阶量}\par\nobreak
分别写两种面积，得到
\[A_h'=\pi a^2(1+\alpha\Delta T)^2,\qquad
A_m'=\pi(b^2-a^2)(1+\alpha\Delta T)^2.\]
令$\varepsilon=\alpha\Delta T$。先扣除初面积，并完整展开平方差，
\[
\begin{aligned}
\Delta A_h&=\pi a^2[(1+\varepsilon)^2-1]
=\pi a^2(2\varepsilon+\varepsilon^2),\\
\Delta A_m&=\pi(b^2-a^2)[(1+\varepsilon)^2-1]
=\pi(b^2-a^2)(2\varepsilon+\varepsilon^2).
\end{aligned}
\]
在$|\varepsilon|\ll1$时舍去二阶项$\varepsilon^2$，才有一阶变化
\[\boxed{\Delta A_h=2\alpha\Delta T\,\pi a^2,\qquad
\Delta A_m=2\alpha\Delta T\,\pi(b^2-a^2)}.\]


\par\medskip\Needspace{4\baselineskip}\noindent\textbf{比较相对变化与绝对增量并核对约束条件}\par\nobreak
孔面积与材料面积的相对变化相同，都是 $2\alpha\Delta T$；它们的绝对增量则由各自初面积决定。若外缘被夹紧，图示整体缩放不再成立，不能据本题结论判断孔的变化。
\end{drillsolution}
\begin{drillsolution}{t1-11}\textbf{隔离高于槽面的液柱并对圆周张力取竖直分量}\par\nobreak
取图\ref{fig:t1-11s}中每根管内高于槽面的液柱为研究对象。将液柱底部压力与自由表面处的均匀大气压部分合并后，均匀$p_0$的合力抵消，剩下的向上作用是接触线处表面张力的竖直合力。对半径为 $r$ 的管，接触线长度为 $2\pi r$，且液体内接触角为 $60^\circ$，故
\[
F_z=\int_0^{2\pi}\sigma\cos60^\circ\,r\,\dd\varphi
=\sigma r\cos60^\circ[\varphi]_0^{2\pi}
=2\pi r\sigma\cos60^\circ.
\]


\par\medskip\Needspace{4\baselineskip}\noindent\textbf{把圆周张力与液柱重力平衡，求细管高度}\par\nobreak
忽略弯月面的体积，液柱重力为 $\rho g\pi r^2h_r$。两者平衡给出所求的高度。\fig{fig:t1-11s}{把两根液柱分别隔离；圆周张力随半径一次变化，单位高度重力随半径平方变化}{\begin{tikzpicture}

\draw[world](0,0)rectangle(.8,2);\fill[blue!8](0,0)rectangle(.8,2);\draw[force](.4,2)--(.4,3.1)node[above]{$\pi r\sigma$};\draw[force](.4,1)--(.4,-.8)node[below]{$\rho g\pi r^2h_r$};\node at(.4,-1.6){半径$r$};
\begin{scope}[xshift=5.5cm]\draw[world](0,0)rectangle(1.6,1);\fill[blue!8](0,0)rectangle(1.6,1);\draw[force](.8,1)--(.8,2.5)node[above]{$2\pi r\sigma$};\draw[force](.8,.5)--(.8,-.8)node[below]{$4\rho g\pi r^2h_{2r}$};\node at(.8,-1.6){半径$2r$};\end{scope}

\end{tikzpicture}}

\[\boxed{h_r=\frac{2\sigma\cos60^\circ}{\rho gr}=\frac{\sigma}{\rho gr}}.\]


\par\medskip\Needspace{4\baselineskip}\noindent\textbf{对粗管重新计周长和面积，求高度比并检查润湿方向}\par\nobreak
把半径换成 $2r$，同样列式，得
\[\boxed{h_{2r}=\frac{\sigma}{2\rho gr},\qquad h_r:h_{2r}=2:1}.\]
图中细管的液面更高，弯月面更弯曲。半径加倍时，张力只增大到2倍，单位高度液柱重力却增大到4倍，故平衡高度减半。接触角从液体一侧量取，$\cos60^\circ>0$，两管都应上升；若把图中弯月面画成不润湿的凸面，便与题给角度和所得正高度不符。
\end{drillsolution}
\begin{drillsolution}{t1-12}\textbf{对整个封闭系统列总物质的量守恒}\par\nobreak
阀门打开后，两室达到同一末温和公共末压。图\ref{fig:t1-12q}中两室始终刚性，总体积为 $3V$；流动只是在内部重新分配分子，总物质的量保持不变。初态分别有
\[n_{1i}=\frac{3p_0V}{RT_0},\qquad n_{2i}=\frac{2p_0V}{RT_0}.\]


\par\medskip\Needspace{4\baselineskip}\noindent\textbf{在公共末温和总体积下求公共末压}\par\nobreak
末态整个系统为 $n_f=p_f(3V)/(RT_0)$。令 $n_f=n_{1i}+n_{2i}$，得
\[\boxed{p_f=\frac{3p_0V+p_0(2V)}{3V}=\frac53p_0}.\]
这相当于按体积加权平均初压，不能将两压强作等权平均。\fig{fig:t1-12s}{阀门两侧的气体重新分配，总量与总体积保持不变}{\begin{tikzpicture}[scale=.8]
\foreach \yy/\pa/\pb/\lab in {0/{3p_0}/{p_0}/初态,-3/{5p_0/3}/{5p_0/3}/末态}{\begin{scope}[yshift=\yy cm]
\draw[thick](0,0)rectangle(2,1.6);\draw[thick](3,0)rectangle(7,1.6);\draw[thick](2,.8)--(3,.8);
\node at(1,1.1){$V,\ \pa$};\node at(5,1.1){$2V,\ \pb$};\node[left]at(-.3,.8){\lab};\node at(1,.4){$T_0$};\node at(5,.4){$T_0$};\end{scope}}
\draw[->,blue,thick](2.15,.8)--(2.85,.8);\node[above]at(2.5,1.8){净流出第一室};\end{tikzpicture}}



\par\medskip\Needspace{4\baselineskip}\noindent\textbf{对第一室单独比较初末物质的量并核对流向}\par\nobreak
图\ref{fig:t1-12s}中第一室末态仍占体积 $V$，所以 $n_{1f}=5p_0V/(3RT_0)$。由第一室净流出的量为
\[\boxed{\Delta n_1=n_{1i}-n_{1f}=\frac{4p_0V}{3RT_0}}.\]
第一室压强由 $3p_0$ 降至 $5p_0/3$，第二室由 $p_0$ 升至 $5p_0/3$，方向与净流动一致。还可核验第二室增加的气体量恰为 $\Delta n_1$，而末态物质的量之比为 $1:2$，与两个容器的体积比相同。
\end{drillsolution}
\begin{drillsolution}{t1-13}\textbf{由自由平衡压强确定接触挡板的温度}\par\nobreak
初态活塞自由平衡。气体向上推力、大气向下压力和活塞重力满足
\[p_*S=p_0S+mg,\qquad p_*=p_0+\frac{mg}{S}.\]



在碰到上挡板前，活塞缓慢上升，负载不变，故气体保持压强 $p_*$。由 $V\propto h$ 和等压关系，气柱高从 $h_0$ 增至 $2h_0$ 时温度为 $2T_0$。题给末温 $3T_0$ 更高，因此后半段必须使用挡板固定体积的条件。\fig{fig:t1-13s}{上挡板向下压住活塞，气体推力须同时平衡外压、重力和挡板反力}{\begin{tikzpicture}

\draw[world](0,.6)--(2,.6);\fill[pattern=north east lines](0,.4)rectangle(2,.6);\draw[force](.4,-.4)--(.4,.35);\node[left]at(.4,-.1){$p_fS$};\draw[force](1,1.6)--(1,.65)node[midway,left]{$p_0S+mg$};\draw[force,blue](1.8,1.6)--(1.8,.65)node[midway,right]{$N$};

\end{tikzpicture}}



\par\medskip\Needspace{4\baselineskip}\noindent\textbf{保持接触后的体积不变，求末压}\par\nobreak
从刚接触态到末态，$V=2Sh_0$ 不变，有
\[\boxed{p_f=p_*\frac{3T_0}{2T_0}=\frac32\left(p_0+\frac{mg}{S}\right)}.\]


\par\medskip\Needspace{4\baselineskip}\noindent\textbf{对末态活塞列受力平衡并检查挡板反力}\par\nobreak
再按图\ref{fig:t1-13s}的末态受力图求支持力。上挡板必须向下压住活塞，故
\[p_fS=p_0S+mg+N,\qquad
\boxed{N=\frac12(p_0S+mg)},\]
方向竖直向下。刚接触时 $N=0$，此后随温度升高而增加；末态 $N>0$，验证了活塞保持接触的假设。若把全过程当作等压，会得到气柱高 $3h_0$，直接越过题图中的上挡板。
\end{drillsolution}

\end{document}
